\documentclass{article}
\usepackage[sglblindworkshop]{neurips_2026}

\usepackage[utf8]{inputenc}
\usepackage[T1]{fontenc}
\usepackage[hidelinks]{hyperref}
\usepackage{url}
\usepackage{booktabs}
\usepackage{amsfonts,amsmath,amssymb,amsthm}
\usepackage{graphicx}
\graphicspath{{figs/}}
\usepackage{microtype}
\usepackage{xcolor}
\usepackage{pifont}

\newtheorem{theorem}{Theorem}
\newtheorem{proposition}[theorem]{Proposition}
\newtheorem{lemma}[theorem]{Lemma}
\newtheorem{conjecture}[theorem]{Conjecture}
\theoremstyle{remark}
\newtheorem{remark}[theorem]{Remark}
\newtheorem{observation}[theorem]{Observation}

\newcommand{\Ex}{\mathbb{E}}
\newcommand{\Ot}{\widetilde{O}}
\newcommand{\etwo}{E_2}
\newcommand{\lamtwo}{\Lambda_2}
\newcommand{\Sset}{\mathcal{S}}
\newcommand{\cm}{\ding{51}}
\newcommand{\xm}{\ding{55}}

\workshoptitle{Second Workshop on ML$\times$OR: Mathematical Foundations and Operational
Integration of Machine Learning for Uncertainty-Aware Decision-Making}

\newcommand{\vhat}{\widehat{\bar v}}
\newcommand{\repourl}{https://github.com/melikabaghi/state-exp3}
\title{Pooling and Drift in Delayed Bandits}
\author{%
  Melika Baghi \\
  Georgia Institute of Technology \\
  \texttt{mbaghi3@gatech.edu}
}

\begin{document}
\maketitle

\begin{abstract}
A system often has to act long before it learns whether the act worked: a recommender sees a click
in seconds and a purchase in days. With $K$ actions and a delay of $d$ rounds, the best rate known for
this setting is $\Ot(\sqrt{(K+d)T})$ over $T$ rounds \citep{esposito2023}, so a longer menu is
always more expensive to learn from. It need not be: if the outcome depends on the action only
through the state it produced, a click or a browse or a cart, then one late outcome informs every
action that could have produced the state observed, and the price is set by how many genuinely
different states the actions produce rather than by how many actions there are. We measure that by
an \emph{effective dimension} $v_t$ between $1$ and the number of states, and prove
$\Ot(\sqrt{(d+1)V\log K})$ for a rotating algorithm and $\Ot(\sqrt{V^-}+\sqrt{dT})$ for the
single-copy one that is actually used, for any budget fixed in advance; merging similar states
lowers the price again, at a bias we make explicit. Even when handed the exact
losses from $d$ rounds ago, no algorithm escapes
$\Omega(\sqrt{d\mathcal E\min\{1+\log J,T/d\}})$, where $J$ counts the drifting directions and
$\mathcal E$ bounds how far the losses move while the learner waits. On generated data the state channel
cuts regret by up to $79$ per cent against action-level weighting, and on the funnel family by
$32$ to $68$ per cent against the minimax-optimal method of \citet{zimmert2020} under the tuning
protocol we report.
\end{abstract}

\section{Introduction}

Delayed outcomes are common in sequential decision-making: a system must act now even though the
outcome it cares about may arrive much later. In recommendation, an item is shown immediately, an
engagement signal such as a click or browsing depth is observed soon after, and conversion or
retention may only be known days or weeks later, so decisions are made before their
economically relevant consequences appear. The central question is when such a signal can reduce
the cost of learning before the final outcome arrives.

The pattern is not special to recommendation: a markdown is set and sell-through is known a season
later, staff are rostered and the shift's outcome lands later, marketing spend commits before
attribution resolves. Here an action is an item that could be shown, a state is what the
user does next, and the delayed outcome is whether they eventually buy. Different items may induce similar distributions over engagement states, so
an outcome observed after one item can carry information about many others. The difficulty scales
not with the number of actions $K$ but with how many genuinely different operational regimes those
actions induce, so thousands of items funnelling users through a few engagement patterns stay easy
to learn from.

The question sits between machine learning and operations research. Exploiting the channel from
actions to states asks what can be inferred from data; the wait before the outcome arrives is a
constraint on the operation, which no estimator removes; and the decision commits while only the
first of the two has been resolved. Bandits with
intermediate observations formalize this: the learner repeatedly picks one action and sees the
consequence of that action alone. \citet{esposito2023} show that the \emph{state-to-loss} map decides the complexity, an
unrestricted one giving only the rate available with no intermediate observations; we fix that map and ask how much the \emph{action-to-state} structure helps.

Intermediate observations face two independent limitations: whether information can be shared
\emph{across actions}, and whether it stays useful \emph{across time}.

Across actions: because similar actions reach similar states, we can \emph{pool through the state},
charging a delayed outcome to every action that could have produced the state observed. What this costs is governed by an
\emph{effective dimension} $v_t$, measuring how distinguishable the action-induced state
distributions are under the current play; it can be far smaller than $K$.

Across time: delay creates a second difficulty even when the learner is handed an exact
description of the past. Let $m_t=c_{t-d}$ be the action-loss vector from $d$ rounds earlier. Its usefulness depends
on how far the losses have moved, which we measure by $\etwo=\sum_t\|c_t-m_t\|_\infty^2$
(Figure~\ref{fig:one}). Our lower bound asks
how much regret this mismatch forces even when $m_t$ is supplied for free.

Adapting to overlap and drift at once remains open
(Section~\ref{sec:conj}).

\begin{figure}[t]
\centering
\includegraphics[width=\textwidth]{fig1.pdf}
\caption{\textbf{An illustration of the two limits on intermediate feedback, on generated data.} \textbf{Left,} what each state predicts about the outcome changes while the map from
actions to states stays fixed, so the losses from $d$ rounds earlier grow increasingly
inaccurate; $\etwo$ of Section~\ref{sec:setting} tracks it. \textbf{Right,} more overlap
between actions allows more sharing, at matched difficulty, $K=40$, $|\Sset|=8$, $T=10000$,
$d=50$, $20$ paired seeds, bars one standard error of the paired difference
(Appendix~\ref{app:num}, study~10).}
\label{fig:one}
\end{figure}

\section{Related work}

Four ingredients separate the neighbouring settings: whether the outcome is delayed, whether a
state is seen, whether the channel producing it is known, and whether the bound is
instance-dependent.

\begin{center}\small
\begin{tabular}{@{}lcccc@{}}
\toprule
 & delayed & state seen & $P$ known & instance-dependent \\
\midrule
\citet{zimmert2020}   & \cm & \xm & \xm & \xm \\
\citet{esposito2023}  & \cm & \cm & \xm & \xm \\
\citet{eldowa2024}    & \xm & \cm & \cm & \cm \\
\citet{baron2025}     & \cm & \xm & \xm & \cm \\
this paper            & \cm & \cm & \cm & \cm \\
\bottomrule
\end{tabular}\\[2pt]
Where this work sits: the four ingredients that separate the neighbouring settings.
\end{center}

The last row is the contribution: no prior work combines all four. \citet{eldowa2024} is closest,
and $v_t-1$ is exactly the $\chi^2$ mutual information they introduce, so that quantity is theirs;
what is new is its role once the outcome is delayed (Appendix~\ref{app:transfer}).
\citet{vernade2020} vary the action-to-state map while fixing the state-to-loss map, where we do
the reverse; \citet{wang2026} and \citet{zhang2026impatient} measure a prediction error as $\etwo$
does, without the delay; and feedback graphs \citep{alon2015,esposito2022,gabbianelli2023} reveal
other actions' losses, where here the state's loss must still be estimated.

\section{Problem formulation}\label{sec:setting}

The question is whether the intermediate state can lower the cost of learning below what $K$ alone
forces. The model that makes it precise is the following.
There are $T$ rounds, $K$ actions, a finite state space $\Sset$, a known action-to-state matrix
$P$, and a fixed known delay $d$; all
logarithms are natural and $\Ot$ hides logarithmic factors. Every round follows the chain
$A_t\xrightarrow{\;P\;}S_t\xrightarrow{\;\theta_t\;}X_t$: the learner draws $A_t$ from a distribution $x_t$,
observes $S_t\sim P(\cdot\mid A_t)$ at once, and receives the delayed outcome only at time $t+d$, encoded as a loss
$X_t\in[0,1]$ so that smaller is better, with
$\Ex[X_t\mid\mathcal F_{t-1},A_t,S_t]=\theta_t(S_t)$, where $\mathcal F_{t-1}$ is everything seen
before round $t$; outcomes are conditionally independent across rounds given the states. The conditional mean depends on $A_t$ only through $S_t$. This
\emph{state-sufficiency} assumption licenses cross-action pooling: an outcome observed at state
$S_t$ may update every action that could have produced that state. If an action affects the outcome
directly, in a way not represented in $S_t$, the argument no longer applies. Here $x_t$ depends only
on $\mathcal F_{t-1}$; $X_r$ may first affect the action
distribution at round $r+d+1$, equivalently the round-$t$ decision uses outcomes only from rounds
$r\le t-d-1$.

The matrix $P\in[0,1]^{K\times|\Sset|}$ is row-stochastic and fixed, and the \emph{state-loss means} $\theta_t\in[0,1]^{|\Sset|}$ are \emph{oblivious}:
chosen in advance, not reacting to the learner. Action losses are $c_t=P\theta_t$ and performance
is the static \emph{pseudo-regret}, the gap to the best single action in hindsight,
$\Ex[R_T]=\Ex[\sum_tc_t(A_t)]-\min_a\sum_tc_t(a)$. \textbf{Oracle for the lower bound.} For Theorem~\ref{thm:lower} only, the learner also receives the \emph{stale action-loss vector} $m_t\triangleq c_{t-d}$ before choosing $A_t$, with $m_t=c_1$ for $t\le d$. This is not a
function of the observables, so a lower bound surviving it is stronger, and
\textsc{State-EXP3} never uses it. We measure how misleading $m_t$ is by $\etwo=\sum_t\|c_t-m_t\|_\infty^2\le T$, which stays small when
errors are small even if frequent. The geometry of $P$ governs sharing across actions and the
evolution of $\theta_t$ governs staleness across time. 

\section{What the theory guarantees}\label{sec:alg}

Two results follow: the number of actions stops setting the price for the algorithm one would
actually run, and no sharing removes the cost of waiting. A rotating variant and one for merged
states are in Appendix~\ref{app:etc}.

An arriving outcome need not update the action played alone. Every action can be credited by how
likely it was to have produced the observed state. Here $P$ is known and $m_t$ is not given. For each state
$s\in\Sset$ let $q_t(s)=\sum_{a=1}^Kx_t(a)P(s\mid a)$, and give each action $a\in\{1,\dots,K\}$ the
estimate $\hat c_t(a)=P(S_t\mid a)X_t/q_t(S_t)$. It is correct on average, and its noise is governed by the
\emph{effective dimension}
$v_t=\sum_sq_t(s)^{-1}\sum_ax_t(a)P(s\mid a)^2\in[1,|\Sset|]$.
\textsc{State-EXP3} runs $m$ copies of \textsc{EXP3} \citep{auer2002} in rotation at learning rate
$\eta$, copy $i\in\{0,\dots,m-1\}$ taking every $m$th round, and adds $\hat c_t$ to a copy's running total once the outcome is usable. A
round costs $O(K)$ (Appendix~\ref{app:pool}, with pseudocode and the $\chi^2$ reading of $v_t$).

\begin{theorem}\label{thm:m1}
Let $P$ be known and $(\theta_t)$ oblivious, and fix in advance any $V^-$ with
$\sum_t(v_t-1)\le V^-$ always.
Then \textup{\textsc{State-EXP3}} with $m=1$ and any $\eta\le1/(e(d+1))$ satisfies
$\Ex[R_T]\le\sqrt{2(3.3V^-+2dT)\log K}+e(d+1)\log K+d$.
\end{theorem}

So $K$ is replaced by $V^-$, which counts how differently the actions behave rather than how many
there are, and the delay enters additively at $\Ot(\sqrt{V^-}+\sqrt{dT})$, for the algorithm every
experiment below runs. Writing $\bar v$ for the largest $v_t$ over all ways of playing,
$V^-=(\bar v-1)T$ is admissible and depends on $P$ alone. Rotating $d+1$ copies and merging
similar states are handled in Appendix~\ref{app:etc}; none of the three adapts to $\etwo$.

The proof turns on one pathwise identity. The denominator of $\hat c_t$ is the state marginal of
the distribution supplying its numerator, so $\langle x_t,\hat c_t\rangle=X_t\le1$ on every path,
whatever $q_t(S_t)$ is. That bounds how far the play distribution moves while an outcome is in
flight, which is what the delayed cross-term needs; centring by $X_t\mathbf 1$ replaces $v_t$ by
$v_t-1$ (Appendix~\ref{app:m1}).

\paragraph{Where this stops.}\label{sec:lower} How much does waiting cost even when the stale
losses are free?
Theorem~\ref{thm:lower} shows a cost from time alone, not a joint overlap-drift bound: the
construction gives each of the $J$ drifting directions a private state, removing useful pooling and
forcing $J\le|\Sset|-1$ (Appendix~\ref{app:proofs}).

\begin{theorem}\label{thm:lower}
Fix $T$, a delay $1\le d\le T/2$, a drift budget $\mathcal E\le T/4$, and
$1\le J\le\min\{K-1,|\Sset|-1\}$. Every algorithm, even one handed the exact losses from $d$ rounds
ago, meets some instance fixed in advance with $\etwo\le\mathcal E$ on which
$\Ex[R_T]=\Omega(\sqrt{d\,\mathcal E\min\{1+\log J,\,T/d\}})$. Without that help, and if $K\le T$,
the best regret any algorithm can guarantee is
$\Omega(\sqrt{KT}+\sqrt{d\,\mathcal E\min\{1+\log J,\,T/d\}})$, up to constants.
\end{theorem}

The first term is the usual cost of exploring. The second is the cost of delay: it grows with the
wait, how far the losses move during it, and the number $J$ of directions the best fixed action can
use. They come from different hard instances, so adding them overstates the truth by at most a
factor of two.

\section{Experiments}\label{sec:num}

A claim about instances is settled where the two counts it separates, how many actions there are
and how many distinct behaviours they have, come apart. The families below are built so they do. With $K=200$ generated items and six levels of engagement the
estimated dimension is $1.97$,
and the state-pooled variant cuts regret against action-level weighting by $79.3$, $63.6$ and
$38.9$ per cent at $d=10,50,200$. That setting is hand-structured, not fitted to data, so it shows
the mechanism, not a deployed system.

It also beats the minimax-optimal delayed-bandit method of \citet{zimmert2020}, tuned over a grid
under the reported tuning protocol.

\begin{center}\small
\begin{tabular}{@{}rrrrrr@{}}
\toprule
$d$ & action-level & tuned \citep{zimmert2020} & \textsc{State-EXP3} & cut & cut \\
\midrule
 10 & 4577 & 2959 &  949 & $79.3\%$ & $67.9\%$ \\
 50 & 4752 & 3557 & 1728 & $63.6\%$ & $51.4\%$ \\
200 & 5307 & 4735 & 3245 & $38.9\%$ & $31.5\%$ \\
\bottomrule
\end{tabular}\\[2pt]
Funnel family, $K=200$, $|\Sset|=6$, $T=20000$, $8$ paired seeds; cuts are against action-level
weighting and against the tuned baseline (study~16).
\end{center}

These runs use the single-copy algorithm, the one Theorem~\ref{thm:m1} covers;
Appendix~\ref{app:num} also reports the rotating variant. A best-state rule leads seven of nine
settings, but where its assumption fails its regret grows linearly and its spread across seeds is
four times ours.

\begin{figure}[t]
\centering
\includegraphics[width=\textwidth]{fig2.pdf}
\caption{\textbf{Where pooling pays.} The gain should follow the effective dimension rather than
the number of actions, and should survive when $P$ is estimated rather than known. Both hold
across the sweep. Relative gain $(R_{\text{action}}-R_{\text{state}})/R_{\text{action}}$, single
copy, $K=40$, $|\Sset|=8$, $T=5000$, $12$ paired seeds; blue is a gain. $\vhat$ estimates
$\sup_xv(x)$ for the raw map by Monte Carlo, which for the coarsened panel is not the dimension of
the coarsened map that the grouping bound charges. The dashed line compares two upper bounds for
the rotating variant, so it is a reference rather than a prediction for the variant plotted.}
\label{fig:phase}
\end{figure}

Figure~\ref{fig:phase} sweeps row overlap against delay and marks where pooling pays, in the three
settings the studies treat separately: $P$ known, $P$ estimated, and states coarsened.

\section{Discussion}\label{sec:conj}

Pooling helps when many actions lead to similar states, and that is checkable before deployment,
since $P$ is known: an estimate of $\bar v$ well below $K$ indicates the state channel is worth
using.

Theorem~\ref{thm:lower} says what no algorithm can do about the second: the uncertainty left by
waiting is irreducible, surviving even when the exact losses from $d$ rounds ago are handed over
free. Two limits remain: $\bar v$ is estimated rather than solved, and the bounds lie on different
axes. 

\begin{thebibliography}{99}
\bibitem[Alon et al.(2015)]{alon2015} N.~Alon, N.~Cesa-Bianchi, O.~Dekel, and T.~Koren. Online learning with feedback graphs: beyond bandits. \emph{COLT}, PMLR 40, 2015.
\bibitem[Auer et al.(2002)]{auer2002} P.~Auer, N.~Cesa-Bianchi, Y.~Freund, and R.~E.~Schapire. The nonstochastic multiarmed bandit problem. \emph{SIAM Journal on Computing}, 32(1):48--77, 2002.
\bibitem[Bar-On and Mansour(2025)]{baron2025} Y.~Bar-On and Y.~Mansour. Non-stochastic bandits with evolving observations. \emph{ALT}, PMLR 272, 2025.
\bibitem[Cesa-Bianchi and Lugosi(2006)]{cesabianchilugosi2006} N.~Cesa-Bianchi and G.~Lugosi. \emph{Prediction, Learning, and Games}. Cambridge University Press, 2006.
\bibitem[Cesa-Bianchi et al.(2019)]{cesabianchi2019} N.~Cesa-Bianchi, C.~Gentile, and Y.~Mansour. Delay and cooperation in nonstochastic bandits. \emph{JMLR}, 20, 2019.
\bibitem[Eldowa et al.(2024)]{eldowa2024} K.~Eldowa, N.~Cesa-Bianchi, A.~M.~Metelli, and M.~Restelli. Information capacity regret bounds for bandits with mediator feedback. \emph{JMLR}, 25(353), 2024.
\bibitem[Esposito et al.(2023)]{esposito2023} E.~Esposito, S.~Masoudian, H.~Qiu, D.~van~der~Hoeven, N.~Cesa-Bianchi, and Y.~Seldin. Delayed bandits: When do intermediate observations help? \emph{ICML}, PMLR 202, 2023.
\bibitem[Esposito et al.(2022)]{esposito2022} E.~Esposito, F.~Fusco, D.~van~der~Hoeven, and N.~Cesa-Bianchi. Learning on the edge: online learning with stochastic feedback graphs. \emph{NeurIPS}, 35, 2022.
\bibitem[Gabbianelli et al.(2023)]{gabbianelli2023} G.~Gabbianelli, G.~Neu, and M.~Papini. Online learning with off-policy feedback. \emph{ALT}, PMLR 201, 2023.
\bibitem[Thune et al.(2019)]{thune2019} T.~S.~Thune, N.~Cesa-Bianchi, and Y.~Seldin. Nonstochastic multiarmed bandits with unrestricted delays. \emph{NeurIPS}, 32, 2019.
\bibitem[Vernade et al.(2020)]{vernade2020} C.~Vernade, A.~Gy\"orgy, and T.~Mann. Non-stationary delayed bandits with intermediate observations. \emph{ICML}, PMLR 119, 2020.
\bibitem[Wang et al.(2026)]{wang2026} S.~Wang, A.~Barik, and V.~Y.~F.~Tan. Bandit convex optimization with gradient prediction adaptivity. \emph{arXiv:2605.22191}, 2026.
\bibitem[Zhang et al.(2026)]{zhang2026impatient} K.~W.~Zhang, T.~Baldwin-McDonald, K.~Ciosek, L.~Maystre, and D.~Russo. Impatient bandits: optimizing for the long-term without delay. \emph{JMLR}, to appear, 2026.
\bibitem[Zimmert and Seldin(2020)]{zimmert2020} J.~Zimmert and Y.~Seldin. An optimal algorithm for adversarial bandits with arbitrary delays. \emph{AISTATS}, PMLR 108, 2020.
\end{thebibliography}

\appendix

\section{The state-pooled estimator}\label{app:pool}

\paragraph{The algorithm.} Inputs are the action-to-state matrix $P$, the delay $d$, the number
of copies $m$, and the rate $\eta>0$. The interleaved setting is $m=d+1$ and the single-copy one
is $m=1$; Theorems~\ref{thm:sexp3} and~\ref{thm:m1} cover them respectively. Round $r$'s outcome becomes usable at round $r+d+1$, which for $m=d+1$ is the next
round belonging to the copy that played round $r$.

\begin{enumerate}\itemsep1pt \parskip0pt \topsep3pt
\item Set $L_i\leftarrow(0,\dots,0)\in\mathbb R^K$ for $i=0,\dots,m-1$, and let $\mathcal P$ be
an empty table of rounds awaiting their outcome.
\item For $t=1,\dots,T$:
\begin{enumerate}\itemsep1pt \parskip0pt \topsep1pt
\item \emph{Deliver.} If $r\triangleq t-d-1\ge1$, read $(i_r,q_r(S_r),S_r,X_r)$ from
$\mathcal P$, remove it, and update only that copy,
$L_{i_r}(a)\leftarrow L_{i_r}(a)+P(S_r\mid a)X_r/q_r(S_r)$ for every $a$. This is the pooled
step, and it touches every action rather than the one played.
\item \emph{Select the copy.} $i\leftarrow t\bmod m$.
\item \emph{Play.} $x_t(a)\propto\exp(-\eta L_i(a))$, draw $A_t\sim x_t$, and observe $S_t$
at once.
\item \emph{Record.} Compute the scalar $q_t(S_t)=\sum_ax_t(a)P(S_t\mid a)$ and store
$(i,q_t(S_t),S_t,\cdot)$ in $\mathcal P$, to be completed when $X_t$ arrives at time $t+d$.
\end{enumerate}
\end{enumerate}

Steps (a), (c) and (d) are each $O(K)$ and there is no optimization subproblem, so a round costs
$O(K)$ once $P$ is stored. The state is held in the $m$ vectors $L_i$, which is $mK$ numbers, the
matrix $P$, which is $K|\Sset|$ numbers, and four scalars for each of the at most $d$ rounds in
flight. Only the scalar $q_r(S_r)$ is needed from round $r$, not the vector $x_r$.

\paragraph{A worked example.} Take three actions and two states, with
\[
P=\begin{pmatrix}0.8&0.2\\0.4&0.6\\0&1\end{pmatrix},
\qquad x_t=(\tfrac12,\tfrac14,\tfrac14),
\]
so $q_t(s_1)=\tfrac12(0.8)+\tfrac14(0.4)+\tfrac14(0)=0.5$ and $q_t(s_2)=0.5$. The third action
cannot produce $s_1$ at all. Suppose the learner draws $a_1$, observes $s_1$, and receives $X_t=1$
at time $t+d$. Then $\hat c_t(a)=P(s_1\mid a)/0.5$, so
$\hat c_t=(1.6,\,0.8,\,0)$ against the action-level $\tilde c_t=(2,0,0)$. One observation updates
all three actions: $a_2$ is charged although it was never played, because it had a $0.4$ chance of
producing the state seen, and $a_3$ is charged nothing, because it could not have produced that
state. Nothing is biased by this. With $\theta_t=(1,0)$ the true losses are $c_t=(0.8,0.4,0)$, and
averaging $\hat c_t$ over $S_t\sim q_t$ returns exactly that. What is bought is the second moment:
$v_t=1.44$ here, against $K=3$ for the action-level estimate, and with $\theta_t=(1,1)$ both are
attained exactly.

\begin{remark}[The $\chi^2$ form]\label{rem:chi2}
Under $x_t$ the pair $(A_t,S_t)$ has law $x_t(a)P(s\mid a)$ with marginals $x_t$ and $q_t$, so
\begin{align*}
\chi^2\bigl(P_{A_t,S_t}\,\big\|\,P_{A_t}\!\otimes\!P_{S_t}\bigr)
&=\sum_{a,s}\frac{\bigl(x_t(a)P(s\mid a)-x_t(a)q_t(s)\bigr)^2}{x_t(a)q_t(s)}\\
&=\sum_{a,s}\frac{x_t(a)P(s\mid a)^2}{q_t(s)}-2+1\ =\ v_t-1,
\end{align*}
the cross term summing to $\sum_{a,s}x_t(a)P(s\mid a)=1$ and the last to
$\sum_{a,s}x_t(a)q_t(s)=1$. So $v_t=1$ exactly when $A_t$ and $S_t$ are independent, which pins the rows of $P$ only on the
support of $x_t$ and therefore means all rows agree whenever $x_t$ has full support, as
exponential weights does. If $P$ is deterministic and every state is reached by some action in
the support of $x_t$, then $v_t$ equals the number of states that support reaches. In particular,
under full-support play and an onto deterministic $P$, $v_t=|\Sset|$. This is an identity, not a bound; it is used for reading $v_t$, and once in the proof of
Theorem~\ref{thm:m1}.
\end{remark}

\begin{proposition}\label{prop:pool}
Fix $t$, let $x_t$ be the play distribution, $q_t(s)=\sum_ax_t(a)P(s\mid a)$ the induced state
distribution, and $\hat c_t(a)=P(S_t\mid a)X_t/q_t(S_t)$, which is well defined because only
states with $q_t(s)>0$ occur. For statements involving $1/x_t(a)$, restrict to actions with
$x_t(a)>0$; \textup{\textsc{State-EXP3}} itself has full support on every action. Then $\hat c_t(a)\ge0$, $\hat c_t(a)\le1/x_t(a)$,
$\Ex[\hat c_t(a)\mid\mathcal F_{t-1}]=c_t(a)$, and
\[
\sum_ax_t(a)\Ex[\hat c_t(a)^2\mid\mathcal F_{t-1}]\ \le\ v_t
:=\sum_{s:q_t(s)>0}\frac{\sum_ax_t(a)P(s\mid a)^2}{q_t(s)},\qquad 1\le v_t\le|\Sset| ,
\]
states of zero probability being omitted throughout, which is the convention $0/0=0$.
The action-level $\tilde c_t(a)=\mathbf 1\{A_t=a\}X_t/x_t(a)$ has the same first three properties
and weighted second moment at most $K$.
\end{proposition}

\begin{proof}
Nonnegativity is immediate from $X_t\ge0$. Given $\mathcal F_{t-1}$ the state $S_t$ has
distribution $q_t$, and $\Ex[X_t\mid\mathcal F_{t-1},S_t=s]=\theta_t(s)$ by the model, since that
conditional mean does not depend on the action, so
$\Ex[\hat c_t(a)\mid\mathcal F_{t-1}]=\sum_sq_t(s)P(s\mid a)\theta_t(s)/q_t(s)=c_t(a)$. For the
range, $q_t(s)\ge x_t(a)P(s\mid a)$ for every $a$, so $P(s\mid a)/q_t(s)\le1/x_t(a)$ and
$X_t\le1$. For the second moment, $X_t^2\le1$ gives
$\Ex[\hat c_t(a)^2\mid\mathcal F_{t-1}]\le\sum_sP(s\mid a)^2/q_t(s)$, and exchanging the sums
gives $v_t$. Each summand of $v_t$ is at most $1$, since $P(s\mid a)\le1$ makes
$\sum_ax_t(a)P(s\mid a)^2\le q_t(s)$, and at least $q_t(s)$, since Jensen applied to
$a\mapsto P(s\mid a)$ under $x_t$ gives $\sum_ax_t(a)P(s\mid a)^2\ge q_t(s)^2$. Summing the two
bounds gives $1\le v_t\le|\Sset|$. The action-level claims are the same computation with the
state replaced by the played action, whose weighted second moment is $\sum_ax_t(a)/x_t(a)=K$.
\end{proof}

The two ends are attained. If every action induces the same state distribution then
$v_t=\sum_sq_t(s)=1$, and if $P$ is deterministic, each action reaching a single state, then
every $P(s\mid a)$ is $0$ or $1$ and each summand is exactly $1$, so $v_t$ counts the states
reached and equals $|\Sset|$ when every state has an action reaching it. Disjoint row supports
give $v_t=K$ rather than $|\Sset|$, since the summands over one action's private states add to
$\sum_sP(s\mid a)=1$, so a sensor with many private states per action resolves few of them. So $v_t$ reads the
overlap between the rows of $P$ rather than their length, and $\bar v=\sup_xv(x)$, the supremum
over play distributions, depends on $P$ alone. An outcome observed at a state updates the estimate
for every action that can reach it, which is why the construction of Section~\ref{sec:lower} gives
each drifting coordinate a private state, forcing $v_t=|\Sset|$ there. The immediately observed
pairs $(A_t,S_t)$ identify $P$ separately, though Theorem~\ref{thm:sexp3} assumes it known.

\section{Proofs: pooling and grouping}\label{app:etc}

Two results are stated here rather than in the body, since Theorem~\ref{thm:m1} covers the
algorithm actually run. The first is the interleaved variant.

\begin{theorem}\label{thm:sexp3}
Let $P$ be known and $(\theta_t)$ oblivious. Pick before the run any fixed number $V$ with
$\sum_tv_t\le V$ always. Then \textup{\textsc{State-EXP3}} with $m=d+1$ and the matching learning
rate satisfies $\Ex[R_T]\le\sqrt{2(d+1)V\log K}$. Writing $\bar v$ for the largest $v_t$ over all
ways of playing, $V=\bar vT\le|\Sset|T$ is one such choice, and it depends on $P$ alone.
\end{theorem}

The second merges states that are not worth telling apart. Group them with a map $g$ from $\Sset$
onto a smaller set $\mathcal Z$ and run the same algorithm on $g(S_t)$, where
$P^g(z\mid a)=\sum_{s:g(s)=z}P(s\mid a)$. Grouping never raises the dimension, but the estimate now
tracks a group average, so the widest spread of losses inside a group,
$\delta_t(g)=\max_z\max_{s,s':g(s)=g(s')=z}|\theta_t(s)-\theta_t(s')|$, is the error introduced.

\begin{theorem}\label{thm:rep}
Under the hypotheses of Theorem~\ref{thm:sexp3}, running on $g$ with $\sum_tv_t(g)\le V(g)$ gives
$\Ex[R_T]\le\sqrt{2(d+1)V(g)\log K}+2\sum_t\delta_t(g)$.
\end{theorem}

Merging need not preserve state-sufficiency, so the grouped estimate is right not for $c_t$ but for
a stand-in loss differing from it by at most $\delta_t(g)$ on every action. Keeping states apart
preserves differences that matter but learns more slowly; merging shares more but pays
$2\sum_t\delta_t(g)$. The best grouping is the true one: with $16$ states from four hidden groups,
sizes $1,2,4,8,16$ give regret $1590$, $831$, $754$, $920$, $1080$, lowest at four, which no
algorithm was told. The map $g$ is fixed in advance; learning it from data remains open.

\textbf{Proof idea.} The pooled estimator is unbiased for every action, while
Proposition~\ref{prop:pool} bounds its play-weighted second moment by $v_t$ rather than $K$. With
$m=d+1$ interleaved copies, the outcome generated on a round belonging to one copy is available
before that copy acts again, so each copy is an ordinary exponential-weights process with
immediate feedback. Summing the $d+1$ potential inequalities gives $(d+1)\log K/\eta$, and the
pooled second moments give $\tfrac\eta2\sum_tv_t$.

\begin{proof}[Proof of Theorem~\ref{thm:sexp3}]
Fix a copy $i$ and let $T_i$ be the number of rounds it owns, so $\sum_iT_i=T$. Round $t$'s
outcome is usable from round $t+d+1$, which with $m=d+1$ is exactly $t+m$, the copy's next round,
so within a copy the feedback is undelayed and $L_i$ carries exactly the estimates of that copy's
earlier rounds. Exponential weights on nonnegative losses give, for any action $a^\star$ and any
$\eta>0$,
\[
\sum_{t\in i}\bigl(\langle x_t,\hat c_t\rangle-\hat c_t(a^\star)\bigr)
\le\frac{\log K}{\eta}+\frac{\eta}{2}\sum_{t\in i}\sum_ax_t(a)\hat c_t(a)^2 ,
\]
by the usual potential telescoping with $e^{-z}\le1-z+z^2/2$, which needs $z\ge0$ and therefore
$\hat c_t\ge0$, supplied by Proposition~\ref{prop:pool}. Each $x_t$ is
$\mathcal F_{t-1}$-measurable, because $L_i$ then holds only rounds $r\le t-m=t-d-1$, so taking
expectations turns the left side into $\sum_{t\in i}(\langle x_t,c_t\rangle-c_t(a^\star))$ by
unbiasedness and the quadratic term into at most $\Ex[\sum_{t\in i}v_t]$, again by
Proposition~\ref{prop:pool}. Summing over the $m=d+1$ copies and using
$\sum_i\min_a\sum_{t\in i}c_t(a)\le\min_a\sum_tc_t(a)$, which is what lets the copies be compared
against one common action, gives $\Ex[R_T]\le(d+1)\log K/\eta+\tfrac\eta2\Ex[\sum_tv_t]$. If
$\sum_tv_t\le V$ surely then $\eta=\sqrt{2(d+1)\log K/V}$ balances the two terms and gives
$\sqrt{2(d+1)V\log K}$, and $V=\bar vT$ is admissible by the definition of $\bar v$.
\end{proof}

The delay enters as a factor because the copies do not share their estimates, each paying
$\log K/\eta$ over $T/(d+1)$ rounds. Sharing them is exactly the $m=1$ algorithm of
the open problem of Section~\ref{sec:conj}.

\begin{lemma}[Coarsening]\label{lem:coarse}
For every play distribution and every $g$, $v_t(g)\le v_t$.
\end{lemma}

\begin{proof}
It suffices to merge two states, the general case following by induction on the merges, and a
merged group of zero probability contributes nothing on either side by the convention above.
Write $P_1,P_2$ for the two columns, $A_j=\sum_ax_t(a)P_j(a)^2$, $B_j=\sum_ax_t(a)P_j(a)$ and
$\lambda=B_1/(B_1+B_2)$. By Cauchy--Schwarz,
$(P_1+P_2)^2\le\lambda^{-1}P_1^2+(1-\lambda)^{-1}P_2^2$ pointwise in $a$, so
$\sum_ax_t(a)(P_1+P_2)^2\le(B_1+B_2)(A_1/B_1+A_2/B_2)$, and dividing by the merged denominator
$B_1+B_2$ shows the merged summand is at most $A_1/B_1+A_2/B_2$, the two it replaces.
\end{proof}

\begin{proof}[Proof of Theorem~\ref{thm:rep}]
Running the algorithm on $g$ is running it on the sensor $(\mathcal Z,P^g)$. The surrogate
losses below depend on $x_t$ and so are predictable rather than oblivious, but the proof of
Theorem~\ref{thm:sexp3} uses obliviousness nowhere, only that each $c^g_t$ is
$\mathcal F_{t-1}$-measurable given the play, so it applies unchanged and bounds regret measured
in the surrogate losses
$c^g_t(a)=\sum_zP^g(z\mid a)\bar\theta_t(z)$, where
$\bar\theta_t(z)=\sum_{s:g(s)=z}q_t(s)\theta_t(s)/q_t^g(z)$ is the group mean the grouped
estimator is unbiased for, by the same computation as in Proposition~\ref{prop:pool} with
$S_t$ replaced by $g(S_t)$. Fix $a$ and $z$. Both $\{q_t(s)/q_t^g(z)\}$ and
$\{P(s\mid a)/P^g(z\mid a)\}$ are probability vectors on $\{s:g(s)=z\}$, so their
$\theta_t$-averages lie in a common interval of length $\delta_t(g)$ and differ by at most
$\delta_t(g)$. Weighting by $P^g(z\mid a)$ and summing over $z$ gives
$|c^g_t(a)-c_t(a)|\le\delta_t(g)$ for every $a$, hence
$\langle x_t,c_t\rangle-c_t(a^\star)\le\langle x_t,c^g_t\rangle-c^g_t(a^\star)+2\delta_t(g)$.
Summing over $t$ adds $2\sum_t\delta_t(g)$, and Lemma~\ref{lem:coarse} makes $V(g)$ no larger than
a bound valid for the raw states.
\end{proof}

\section{Relation to other inaccuracy measures}\label{app:transfer}

Write $\lamtwo=\sum_t\min\{1,\|c_t-m_t\|_2^2\}$ for the inaccuracy measure of \citet{baron2025}.
Unlike $\etwo$, which charges only the worst action in a round, $\lamtwo$ grows with how many
actions are predicted wrongly, and since $\|v\|_\infty\le\|v\|_2\le\sqrt J\|v\|_\infty$ for a
vector $v$ with $J$ nonzero coordinates,
\[
\etwo\ \le\ \lamtwo\ \le\ J\,\etwo,
\]
with both ends attained: the left when a single action is mispredicted, the right when all $J$ are
mispredicted equally. So a bound in $\etwo$ is never weaker and can be $J$ times stronger, which is
why Theorem~\ref{thm:lower} is stated in $\etwo$. The two measures agree exactly when the
prediction error is concentrated on one action.

\section{Proofs: the lower bound}\label{app:proofs}

The measure $\etwo$ is controlled by the state-loss variation $W$, and no converse bound
holds in general.

\begin{lemma}[Delay window]\label{lem:window}
$\etwo\le d^2W$ for $W=\sum_t\|\theta_t-\theta_{t-1}\|_\infty^2$. The ratio $\etwo/(d^2W)$ tends to one when $\theta$ drifts monotonically in one coordinate at a
constant rate and $T/d\to\infty$, the first
$d$ rounds contributing $\sum_{j<d}j^2$ rather than $d^2$ each under the convention $m_t=c_1$.
\end{lemma}

In words, accumulated squared prediction error is at most the state-loss variation inflated by
the square of the delay, and no better bound in $W$ alone is available.

\begin{proof}[Proof of Lemma~\ref{lem:window}]
Throughout write $\Delta_j=\theta_j-\theta_{j-1}$, set $\theta_0=\theta_1$, so $\Delta_1=0$, and read
$c_r$ as $c_1$ for $r\le0$, matching the convention $m_t=c_1$ for $t\le d$ of
Section~\ref{sec:setting}. Then $c_t-c_{t-d}=P\sum_{j=t-d+1}^t\Delta_j$. Because the
rows of $P$ are probability vectors, averaging cannot increase the supremum norm, so
$\|c_t-c_{t-d}\|_\infty\le\sum_j\|\Delta_j\|_\infty$. By Cauchy--Schwarz applied to the $d$
summands, $\|c_t-c_{t-d}\|_\infty^2\le d\sum_{j=t-d+1}^t\|\Delta_j\|_\infty^2$. Summing over $t$
and using that each increment belongs to at most $d$ windows gives $\etwo\le d^2W$. The ratio $\etwo/(d^2W)$ approaches one when every increment has the same magnitude, sign and
active coordinate, and some action reaches that coordinate deterministically, since then no
cancellation occurs within a full window and $P$ passes the increment through undiminished.
Without the second condition the drift can lie in the kernel of $P$ and the ratio can be zero.
\end{proof}

\begin{lemma}[Information isolation]\label{lem:iso}
If $h\le d$ then $A_t$ is independent of the signs governing its own block, and this remains true
when the learner is provided with the exact stale action-loss vector $m_t$.
\end{lemma}

\begin{proof}[Proof of Lemma~\ref{lem:iso}]
For $t$ in block $b$, every realized outcome usable when choosing $A_t$ comes from a round
$r\le t-d-1$. Since $t\le bh$ and $h\le d$, we have $r\le bh-d-1<(b-1)h$, so every usable realized
outcome depends only on signs from blocks strictly before $b$, whose joint law depends only on
$\sigma_2,\dots,\sigma_{b-1}$. Hence by conditional
independence $A_t$ is independent of block $b$'s signs. For the oracle, take any $h\le d$. If
$t>d$ then $t-d\le bh-d\le(b-1)h$ by the same computation, so $m_t=c_{t-d}$ is measurable with
respect to the signs of blocks strictly before $b$. If $t\le d$ then $m_t=c_1=\tfrac12\mathbf 1$,
which is deterministic because block~$1$ is neutral. In both cases $m_t$ is independent of block
$b$'s signs, so adjoining it to the learner's information leaves the conclusion intact.
\end{proof}

The construction in full. Take states $s_0,\dots,s_q$, a stationary deterministic map sending
$a_j$ to $s_j$ for $j\le J$ and every other action to $s_0$, and independent Rademacher signs
$\sigma_b^{(j)}$ for $2\le b\le B=\lfloor T/h\rfloor$. Set $\theta_t(s_0)=\tfrac12$ throughout,
$\theta_t\equiv\tfrac12$ on block~$1$, and $\theta_t(s_j)=\tfrac12-\varepsilon\sigma_b^{(j)}$ on
each later block $b$. The neutral first block makes $c_1=\tfrac12\mathbf 1$, so the
convention $m_t=c_1$ for $t\le d$ of Section~\ref{sec:setting} supplies a constant over exactly
the rounds that precede any usable feedback, which is what Lemma~\ref{lem:iso} needs at $b=1$. Rounds after the last complete block carry no drift, costing a constant factor.

\begin{lemma}[Rademacher lower tail]\label{lem:rade}
Let $S_B$ be a sum of $B\ge32$ independent Rademacher signs. For every $1\le u\le\sqrt B/32$,
$\Pr(S_B\ge u\sqrt B)\ge\tfrac u4e^{-64u^2}$.
\end{lemma}

\begin{proof}
Write $S_B=2X-B$ with $X\sim\mathrm{Bin}(B,\tfrac12)$, so that
$\{S_B\ge u\sqrt B\}=\{X\ge B/2+u\sqrt B/2\}$. Passing to $X$ removes the lattice gap, since $X$
takes every integer value in $[0,B]$ whereas $S_B$ takes only those of the parity of $B$. Put
$n=\lfloor B/2\rfloor$ and $t=\lceil u\sqrt B/2\rceil+1$, so $X\ge n+t$ implies the event, and
$t\le u\sqrt B$ because $u\sqrt B\ge\sqrt{32}>4$.

For $j\ge1$ the ratio of consecutive binomial probabilities is
$\Pr(X=n+j)/\Pr(X=n+j-1)=(B-n-j+1)/(n+j)$. Using $(B-1)/2\le n\le B/2$, its distance from one is
at most $(2j-1)/(n+j)\le(4j-2)/(B-1)\le8j/B$ for $B\ge2$. Restricting to $j\le B/16$ keeps
$8j/B\le\tfrac12$, where $\log(1-x)\ge-2x$ holds, so summing over $i\le j$ gives
$\Pr(X=n+j)\ge\Pr(X=n)e^{-16j(j+1)/(2B)}\ge\Pr(X=n)e^{-16j^2/B}$. The central term satisfies
$\Pr(X=n)\ge1/(2\sqrt B)$ for every $B\ge2$.

Sum over the $t$ integers $j\in\{t,\dots,2t-1\}$, all of which are at most $2t$ and hence at most
$B/16$ because $u\le\sqrt B/32$ forces $2t\le2u\sqrt B\le B/16$. Each term is at least
$e^{-64t^2/B}/(2\sqrt B)$, and $t\le u\sqrt B$ makes the exponent at most $64u^2$, so
$\Pr(X\ge n+t)\ge t\,e^{-64u^2}/(2\sqrt B)\ge\tfrac u4e^{-64u^2}$ using $t\ge u\sqrt B/2$.
\end{proof}

Three regimes give the maximal inequality the lower bound needs. Each $S_B^{(j)}$ is sub-Gaussian
with variance proxy $B$, so $\Ex\max_j(S_B^{(j)})^{+}\le\sqrt{2B\log J}+\sqrt B$ throughout. For
the matching lower bound, first suppose $\log J\le B/8$ and $\log J\ge128$. Taking
$u=\sqrt{\log J/128}$, which then lies in $[1,\sqrt B/32]$, makes $e^{-64u^2}=J^{-1/2}$, so
$\Pr(S_B^{(j)}\ge u\sqrt B)\ge\tfrac u4J^{-1/2}$ and independence across $j$ gives
$\Pr(\max_jS_B^{(j)}<u\sqrt B)\le e^{-u\sqrt J/4}$, at most $\tfrac12$ beyond an absolute $J$.
Hence $\Ex\max_j(S_B^{(j)})^{+}\ge\tfrac12u\sqrt B=\Omega(\sqrt{B\log J})$. Second, if
$\log J>B/8$ and $B\ge1024$, take $u=\sqrt B/32$, so $u\sqrt B=B/32$ and $e^{-64u^2}=e^{-B/16}$,
and $J\ge e^{B/8}$ makes $J\,e^{-B/16}$ diverge, giving
$\Ex\max_j(S_B^{(j)})^{+}=\Omega(B)$, which is $\Omega(\sqrt{B\min\{1+\log J,B\}})$ in that
regime. Third, when $\log J<128$ or $B<1024$ the quantity $\sqrt{B\min\{1+\log J,B\}}$ is
$O(\sqrt B)$, and a single walk already gives
$\Ex(S_B^{(1)})^{+}=\tfrac12\Ex|S_B|\ge\tfrac12(\Ex S_B^2)^{3/2}(\Ex S_B^4)^{-1/2}\ge\tfrac12\sqrt{B/3}$
by H\"older, since $\Ex S_B^2=B$ and $\Ex S_B^4\le3B^2$. Combining the three regimes,
$\Ex\max_j(S_B^{(j)})^{+}=\Theta(\sqrt{B\min\{1+\log J,B\}})$. This is the route by which maxima
of independent random walks fix the minimax rate for prediction with expert advice
\citep{cesabianchilugosi2006}.

\textbf{Proof idea.} Divide time into blocks of length $d$ and give $J$ actions private states
whose losses are perturbed by independent Rademacher signs within each block. Every observation
available when an action is chosen comes from an earlier block, and the supplied $m_t=c_{t-d}$
depends only on earlier signs, so the learner cannot predict the signs governing its current
block. Its expected loss is therefore $T/2$, while the best fixed action in hindsight benefits
from the maximum of $J$ independent random walks. Choosing the amplitude to make $\etwo\le\mathcal E$
gives the stated scale.

\begin{proof}[Proof of Theorem~\ref{thm:lower}]
Take $h=d$ and $\varepsilon=\tfrac12\sqrt{\mathcal E/T}$. By Lemma~\ref{lem:iso} every algorithm's
expected loss equals $T/2$, since its action is independent of the current block's signs and all states have mean
$\tfrac12$ marginally. Its regret is therefore exactly the comparator's
fluctuation advantage,
$\varepsilon d\,\Ex\bigl[\max_{j\le J}\bigl(\sum_{2\le b\le B}\sigma_b^{(j)}\bigr)^{+}\bigr]$ with
$B=\lfloor T/d\rfloor$, over the $B-1$ randomized blocks. Per round $\|c_t-m_t\|_\infty\le2\varepsilon$, because the supremum norm charges only the largest
coordinate and each coordinate moves by at most $2\varepsilon$ at a block boundary; the gap is
$\varepsilon$ across the first randomized block, where $m_t$ is still the neutral $c_1$, and lies
in $\{0,2\varepsilon\}$ thereafter. Hence $\etwo\le4\varepsilon^2T=\mathcal E$ holds
deterministically. The
budget must bind on the realized sequence, since a budget holding only in expectation would not
constrain the instance selected by Yao's principle. By Lemma~\ref{lem:rade} and the discussion following it, the expected maximum of the positive
parts of $J$ independent Rademacher sums of length $B-1$ is
$\Theta(\sqrt{(B-1)\min\{1+\log J,B-1\}})$, which covers both regimes and the case $J=1$. The
hypothesis $d\le T/2$ gives $B\ge2$ and hence $B-1\ge B/2$, so replacing $B-1$ by $B$ costs a
factor at most $\sqrt2$ in each branch of the minimum, including the saturated branch where the
minimum is $B-1$ rather than $1+\log J$. Substituting $B=\lfloor T/d\rfloor$ and $\varepsilon$
therefore gives $\Theta(\sqrt{d\mathcal E\min\{1+\log J,T/d\}})$. Rounds after the last complete block
carry no drift and change the constant only. When the oracle is withheld, the $\sqrt{KT}$ term is \citet[Thm.~5.1]{esposito2023}, whose
construction is stationary and therefore lies in the present model class.
\end{proof}

\begin{proposition}\label{prop:tight}
Suppose $X_t=c_t(A_t)$. On instances whose drift is carried by a single action, $\lamtwo=\etwo$,
so the bound of \citet{baron2025} reads $\Ot(\sqrt{KT+d\,\etwo})$ and its drift contribution
$\sqrt{d\,\etwo}$ matches Theorem~\ref{thm:lower} up to logarithmic factors.
\end{proposition}

In words, single-action drift makes the Euclidean and supremum measures agree, so on that family
the drift term of the published upper bound is unimprovable in order. Its $\sqrt{KT}$ term is not
matched because the oracle supplies an entire stale loss vector, so the oracle-assisted problem
need not retain bandit exploration complexity. Neither result provides an $\etwo$-adaptive upper
bound under noisy outcomes.

\begin{proof}[Proof of Proposition~\ref{prop:tight}]
If only one coordinate of $c_t-m_t$ is nonzero then $\|c_t-m_t\|_2=\|c_t-m_t\|_\infty$. On the
family of Section~\ref{sec:lower} that coordinate has magnitude at most $2\varepsilon\le1$, so
the truncation at one in $\lamtwo$ is inactive and $\lamtwo=\etwo$. The bound of \citet{baron2025} then reads $\Ot(\sqrt{KT+d\etwo})$, whose $\sqrt{d\etwo}$
contribution Theorem~\ref{thm:lower} matches up to the logarithmic factor at $J=1$. Its
$\sqrt{KT}$ contribution is not matched.
\end{proof}

\section{Numerical detail}\label{app:num}

Sixteen studies were run and the seven bearing on the claim are reported here, each with an
observation, an interpretation and a limitation. They keep their original numbers, which index the
released scripts, so the numbering is not contiguous; the code for all sixteen is available.
Studies 1 to 13 draw $P$ from a Dirichlet and 14 to 16 do not. \textsc{State-EXP3} is beaten in
four of six cells of study~16 when the baseline alone is grid-tuned, and study~10 is a control for
study~6 that reverses its conclusion. All
policies run on the
same environment parameters and common random-number streams, so a seed fixes everything except
which estimator is formed, and differences are taken within a seed before averaging. Round $r$'s outcome is consumed before round $r+d+1$, matching
Section~\ref{sec:setting}. Reported $\pm$ is one standard error of the paired difference. Throughout this appendix,
\emph{sample-path pseudo-regret} means $\sum_tc_t(A_t)-\min_a\sum_tc_t(a)$, evaluated on one
realization of the learner and environment randomness. Averaging it over seeds estimates the
expected pseudo-regret that the theorems bound. We use that name rather than ``realized regret'' to
distinguish it from regret computed from the sampled outcomes $X_t$, which is not what these
studies report.

\textbf{Study~1. The bound holds and pooling helps.} With $|\Sset|=4$, $T=10000$, $d=50$,
drift band $0.05$ and $100$ seeds:

\begin{center}\small
\begin{tabular}{rrrrrr}
\toprule
$K$ & state, $m=d+1$ & state, $m=1$ & action, $m=1$ & paired gain & state better in \\
\midrule
 8 & 455.1 & 368.2 & 387.6 & $19.3\pm3.4$ &  73/100 \\
20 & 618.3 & 493.7 & 547.5 & $53.8\pm5.6$ &  87/100 \\
40 & 734.3 & 574.5 & 696.1 & $121.5\pm9.4$ &  99/100 \\
80 & 780.4 & 614.8 & 784.0 & $169.3\pm9.9$ & 100/100 \\
\bottomrule
\end{tabular}\\[2pt]
Study 1, $|\Sset|=4$, $T=10000$, $d=50$, 100 paired seeds; $\pm$ is one standard error.
\end{center}

Observed. For the interleaved variant, sample-path pseudo-regret stayed below the bound
$\sqrt{2(d+1)|\Sset|T\log K}$ in all $100$ seeds at every $K$, that bound running from $2913$ to
$4228$ against realized $455$ to $780$. The single-copy variant beat action-level weighting at every
$K$, by a margin growing from $19$ to $169$ as $K$ rose tenfold with $|\Sset|$ fixed, and beat the
interleaved variant throughout. Sweeping the delay at $K=40$ gave state-pooled $387$, $569$
and $731$ against action-level $653$, $699$ and $773$ at $d=10$, $50$ and $200$.

Interpretation. The margin growing in $K$ at fixed $|\Sset|$ is the prediction that separates the
two estimators, so the improvement is attributable to the state channel the effective dimension attributes to it.
Interleaving looks like a price of the analysis rather than of the estimator. The narrowing of the
gap as the delay grows is what an additive delay term predicts and a multiplicative one does not,
which is the diagnostic evidence behind the open problem of Section~\ref{sec:conj}.

Limitation. This is one synthetic family with a fixed drift band, and a simulation cannot
establish that a bound holds, only that sample-path pseudo-regret stayed below it in the runs
performed.

\textbf{Study~6. The effective dimension, not $|\Sset|$, governs the gain.} Holding
$|\Sset|=8$, $K=40$, $T=5000$, $d=20$ and $24$ seeds, and varying only the overlap between the
rows of $P$, the measured mean of $v_t$ moves while $|\Sset|$ does not:

\begin{center}\small
\begin{tabular}{rrrrr}
\toprule
row overlap & mean $v_t$ & state-pooled & action-level & paired gain \\
\midrule
0.00 & 2.719 & 392.7 & 599.3 & $206.6\pm20.0$ \\
0.25 & 2.013 & 361.1 & 528.1 & $167.0\pm21.3$ \\
0.50 & 1.532 & 306.5 & 406.5 & $100.0\pm10.5$ \\
0.75 & 1.168 & 199.9 & 229.2 & $29.3\pm3.2$ \\
0.95 & 1.009 &  49.2 &  50.8 & $1.6\pm0.5$ \\
\bottomrule
\end{tabular}
\end{center}

Observed. Both the regret and the advantage over action-level weighting fall monotonically with
$v_t$, from a paired gain of $206.6$ at $v_t=2.719$ to $1.6$ at $v_t=1.009$. Nothing about the raw
state count changes across the sweep.

Interpretation. $|\Sset|$ is constant and cannot explain the pattern, while $v_t$ moves with it,
which is the comparison Theorem~\ref{thm:sexp3} predicts.

Limitation. The two estimators do not become the same object at $v_t\approx1$. Rather the rows of
$P$ agree, so every action induces the same loss and both regrets collapse, making that endpoint
uninformative rather than favorable. One parameter of one map family was varied, so the sweep is
consistent with the effective dimension governing the gain rather than establishing it.

\begin{figure}[t]
\centering
\includegraphics[width=\textwidth]{fig3.pdf}
\caption{Study~8, the drift channel. \textbf{Left,} measured regret against the scale
Theorem~\ref{thm:lower} predicts, over twelve cells in which $d$ varies $20$-fold, $\etwo$
$100$-fold and $J$ $16$-fold. The dashed line has slope one. \textbf{Right,} the same points
divided by that scale, grouped by which parameter the cell sweeps, with the shaded band spanning
the observed range. \textsc{State-EXP3} appears here only as one more algorithm on the hard
family, where it has no advantage.}
\label{fig:drift}
\end{figure}

\textbf{Study~8. The drift bound describes the right scaling.} The construction of
Theorem~\ref{thm:lower} at $K=40$, $T=8000$, $h=d$ and $12$ paired seeds, sweeping $d$ from $10$
to $200$, the amplitude $\varepsilon$ from $0.02$ to $0.20$ so that $\etwo$ runs from $12$ to
$1193$, and the drifting coordinates $J$ from $1$ to $16$. We run four policies: uniform, action-level \textsc{EXP3}, \textsc{State-EXP3},
and the greedy policy on the exact stale vector $m_t$, the last because
Theorem~\ref{thm:lower} claims to survive it. $\etwo$ is measured on the realized instance rather
than predicted, so the normalization uses the quantity the theorem uses.

Observed. The predicted scale $\sqrt{d\,\etwo\min\{1+\log J,T/d\}}$ ranges from $38$ to $377$
across the grid and the measured regret tracks it (Figure~\ref{fig:drift}). Normalized, the four
policies read $0.230$ to $0.429$, $0.228$ to $0.424$, $0.238$ to $0.417$ and $0.237$ to $0.437$,
a spread below $1.9$ against a tenfold range in the predictor. The oracle-assisted policy has the
largest mean ratio, $0.345$ against $0.311$, $0.314$ and $0.310$ for the three that never see
$m_t$.

Interpretation. The three arguments enter the predicted scale differently and the collapse holds
in all three sweeps, so it is the form of the bound rather than one fitted constant that tracks
the data. That greedy play on the exact $m_t$ sits at the top of the band rather than below it is
the concrete content of the claim that the bound survives the oracle.

Limitation. Every cell has $T/d>1+\log J$, so only the $1+\log J$ branch of the minimum is
exercised and the $T/d$ saturation branch, which needs $d$ comparable to $T$, is not probed. A
lower bound quantifies over all algorithms and four policies cannot establish that. What these
runs check is the scaling of the construction.

\textbf{Code and environment.} The code producing every reported number and figure is at
\expandafter\url\expandafter{\repourl}. It pins the software environment used for the reported
numbers, python 3.14.3 with numpy 2.4.2 and matplotlib 3.10.8, states the seed count for each
study, lists one command per numbered study, and commits the stored outputs each figure reads so
that every figure rebuilds without rerunning an experiment.

\begin{figure}[t]
\centering
\includegraphics[width=\textwidth]{fig5.pdf}
\caption{\textbf{Left,} study~10. The same overlap sweep with and without holding the task
difficulty fixed. Matching the uniform-play regret reverses the trend, supporting the interpretation that the
decline in Study~6 is driven by task-difficulty confounding rather than by the pooling mechanism
itself. Bars are one standard error of the paired
difference. \textbf{Right,} regret saved against action-level weighting over the smoothing
pseudocount and the sensor balance, $10$ paired seeds per cell, from the plug-in sensitivity study
whose script is released with the rest.}
\label{fig:matched}
\end{figure}

\textbf{Study~10. Overlap drives the gain, not easier tasks.}
Study~6 sweeps the overlap and finds the advantage falling as $\vhat$ falls, which read at face
value contradicts Theorem~\ref{thm:sexp3}, whose bound improves as $\bar v$ falls. That sweep is
confounded, since raising the overlap also flattens $c_t=P\theta_t$ across actions, collapsing
the comparator gap so that both regrets vanish together. This study holds $K=40$, $|\Sset|=8$,
$T=10000$, $d=50$ and the difficulty fixed, and varies only $\vhat$. Difficulty is the
uniform-play regret $\sum_t\mathrm{mean}_ac_t(a)-\min_a\sum_tc_t(a)$, a property of the instance
with no algorithm in it, and the amplitude of $\theta$ is bisected in every cell to hit a common
target. The map keeps one contrast direction alive at every overlap and $\theta$ is aligned with
it, which is what lets the gap survive as the rows agree.

Observed. Across $20$ paired seeds the relative gain rises from $11.7$ per cent at
$\vhat=3.380$ to $19.3$ per cent at $\vhat=1.120$, with absolute gains $85.5\pm8.9$ to
$106.2\pm7.9$ and a correlation of $-0.929$ between $\vhat$ and the gain. Study~6's unmatched
sweep falls from $34.5$ to $3.1$ per cent over a comparable range (Figure~\ref{fig:matched}).

Interpretation. The two sweeps disagree in sign, and the matched-difficulty sweep agrees with the overlap
dependence suggested by Theorem~\ref{thm:sexp3}, although this experiment uses the single-copy
variant rather than the interleaved one that theorem analyses. Overlap is what pooling exploits, and study~6
measures overlap confounded with an instance becoming trivial.

Limitation. At the two highest overlaps some seeds cannot reach the common target and the
realized difficulty is $5$ to $13$ per cent below it, so the matching is close rather than exact.
Since a smaller gap should if anything shrink the advantage, the measured rise is conservative in
that direction. The matched gains are also smaller in magnitude than the unmatched ones, which is
what fixing the difficulty costs.

\textbf{Study~10b. The matched-difficulty trend survives a second configuration.} Study~10 is a single
design point, so the same procedure was re-run at $K=20$, $|\Sset|=12$, $T=8000$, $d=20$ and $20$
paired seeds, with the common target lowered to $500$ so that it is reachable at every overlap.
Nothing else changes.

\begin{center}\small
\begin{tabular}{@{}rrrrrr@{}}
\toprule
overlap & $\vhat$ & $R_{\text{unif}}$ & state & action & paired gain \\
\midrule
0.00 & 3.440 & 500.0 & 364.7 & 407.2 & $42.5\pm8.2$ \\
0.35 & 2.225 & 500.0 & 346.8 & 401.3 & $54.6\pm8.3$ \\
0.65 & 1.518 & 500.0 & 309.3 & 372.4 & $63.2\pm6.4$ \\
0.85 & 1.240 & 493.6 & 271.8 & 346.0 & $74.3\pm7.2$ \\
0.95 & 1.191 & 490.1 & 256.9 & 327.6 & $70.7\pm7.8$ \\
\bottomrule
\end{tabular}
\end{center}

Observed. The relative gain rises from $10.4$ to $21.6$ per cent as $\vhat$ falls from $3.440$ to
$1.191$, and the correlation between $\vhat$ and the absolute gain is $-0.972$ against study~10's
$-0.929$. Realized difficulty spans $490.1$ to $500.0$, a $1.020$ times spread, so the matching is
tighter here than in study~10.

Interpretation. The direction study~10 reports is not an artifact of its single configuration. Two
independent design points, differing in $K$, $|\Sset|$, $T$ and $d$, give the same sign and a
comparable magnitude.

Limitation. Both configurations use the same construction and the same matching routine, so this
tests configuration sensitivity and not construction sensitivity. Both run the practical $m=1$
variant. An earlier attempt at this replication silently failed because
\texttt{match\_scale} binds its target as a default argument, so the bisection saturated and the
difficulty was never matched; the run reported here passes the target explicitly.

\begin{figure}[t]
\centering
\includegraphics[width=\textwidth]{fig7.pdf}
\caption{Study~14, a recommendation funnel. \textbf{Left,} the effective dimension stays near two
while the catalogue grows thirty-twofold, so the gap between $K$ and $\vhat$ comes from the
catalogue having categories, not from a chosen $K$. \textbf{Right,} the three arms at $T=20000$
and $8$ paired seeds, labelled by whether the regime condition of Section~\ref{sec:alg} holds.
Bars are one standard error. No real data is used anywhere in this study.}
\label{fig:funnel}
\end{figure}

\textbf{Study~14. Many actions resolved by few states is a natural regime.} Every study so far draws
$P$ from a Dirichlet, which is the assumption under test rather than evidence for it. This one
builds an environment to the shape of a recommendation funnel instead. Actions are $K=200$
catalogue items, states are $|\Sset|=6$ ordered engagement depths, and the outcome is a delayed
conversion. Three structural facts are imposed because funnels have them, namely that items fall
into a modest number of categories sharing an engagement profile, that item quality is
heavy-tailed, and that the loss falls with engagement depth and drifts seasonally rather than
adversarially. The generated catalogue puts $74$ per cent of its mass on no engagement and $1.4$
per cent on the deepest.

Observed. The effective dimension is $\vhat=1.97$ against $K=200$, a hundredfold gap, and it
stays between $1.41$ and $1.81$ as the catalogue grows from $25$ to $800$ items
(Figure~\ref{fig:funnel}). Against action-level weighting the state-pooled variant saves $79.3$,
$63.6$ and $38.9$ per cent at $d=10$, $50$ and $200$, the largest margins anywhere in this
appendix. The plug-in matches known $P$ at every delay, the paired difference never reaching two
standard errors. The regime condition $\bar v(d+1)\log K\le K+d$ holds at $d=10$ and fails at
$d=50$ and $d=200$.

Interpretation. The separation between $K$ and $\vhat$ arrives from catalogue structure rather
than from a tuned parameter, and it does not close as the catalogue grows, which is the claim the
Dirichlet families cannot make. That the condition Theorem~\ref{thm:sexp3} needs fails at the
delays a funnel has, while the single-copy variant still saves most of the regret, is the same
message as study~1 in the setting that motivates the paper. It is evidence for
the open problem of Section~\ref{sec:conj} rather than for the theorem.

Limitation. \emph{No real data is used.} There is no interaction log in this environment, nothing
here is fitted to one, and this is not a semi-synthetic benchmark in the usual sense of that
phrase. The structure is imposed by hand from qualitative properties, so it shows that the regime
is describable and self-consistent, not that any deployed system occupies it. The arms run $m=1$,
covered by Theorem~\ref{thm:m1}, and the environment has no adversarial component, so it exercises the
pooling channel and not the drift channel.

\textbf{Study~16. Pooling still helps against the strongest published algorithm.} Every
comparison so far is against action-level weighting, \textsc{EXP3-IX}, naive sharing, or the
best-state rule. None is the strongest published method for delayed feedback: that is the hybrid
Tsallis-plus-entropy \textsc{FTRL} of \citet{zimmert2020}, whose delay term is additive rather
than multiplicative. We run it on both headline families, with the same environments, seeds and
delay convention, under three tunings.

\begin{center}\small
\begin{tabular}{@{}lrrrr@{}}
\toprule
family, cell & \textsc{State-EXP3} & \citeauthor{zimmert2020} & same, grid-best & \textsc{State-EXP3}, grid-best \\
\midrule
funnel, $d=10$   &  949.0 & 2959.4 & 2413.8 & 27.0 \\
funnel, $d=50$   & 1728.1 & 3557.0 & 2556.4 & 95.9 \\
funnel, $d=200$  & 3244.8 & 4734.7 & 2807.7 & 199.7 \\
overlap $0.00$   &  661.5 &  731.8 &  610.4 & 202.4 \\
overlap $0.65$   &  584.1 &  673.1 &  544.0 &  51.2 \\
overlap $0.95$   &  461.4 &  540.2 &  419.4 &  34.3 \\
\bottomrule
\end{tabular}\\[2pt]
Study 16. Regret under three tunings; grid-best columns are chosen by reading these benchmarks.
\end{center}

Observed. With \textsc{State-EXP3} at the single tuning used everywhere else here and
\citet{zimmert2020} at the best of a six-point grid around its own value, \textsc{State-EXP3} wins
every cell, cutting regret by $31.5$ to $67.9$ per cent on the funnel and $9.6$ to $14.6$ per cent
on the matched sweep. With both at their grid-best, over seven multipliers for
\citet{zimmert2020} and six for \textsc{State-EXP3}, \textsc{State-EXP3} wins every cell again and
by more, $66.8$ to $98.9$ per cent. With \citet{zimmert2020} at its grid-best and
\textsc{State-EXP3} left at its paper tuning, the third column against the first,
\citet{zimmert2020} wins the funnel at $d=200$ and all three matched cells, four of six.

Interpretation. No column is a symmetric protocol. The two grids are finite and unequal, and
\textsc{State-EXP3}'s grid-best sits at its top multiplier in every cell, so its own minimum may
lie outside the range searched; the truncation runs against \textsc{State-EXP3}. Where it is not
held to the worse tuning of the two, the state channel helps against the strongest published
method, which is expected, since that method is action-level by construction and cannot read the
state at all. Grid tuning helps \textsc{State-EXP3} far more than it helps \citet{zimmert2020},
which is consistent with pooling cutting estimator variance and lower variance admitting a larger
learning rate.

Limitation. The third column against the first is a real loss and is reported as one. These are
our own synthetic families, and a method built for adversarial worst cases is not shown at its
best on benign instances. Every grid-best multiplier is chosen by reading these benchmarks, so
none of the last three columns is a procedure a practitioner could run; only the first is. Eight
paired seeds on the funnel and ten on the matched sweep.

Taken together these studies cover both channels. Studies 1, 6, 10 and 14 probe the state channel
and are consistent with the effective dimension, rather than $|\Sset|$ or $K$, governing what
pooling buys; study~10 had to control for task difficulty before study~6's sweep
pointed that way. Study~14 is the only one whose matrix is not drawn from a Dirichlet and reports
the largest margins here, on a structure imposed by hand rather than measured. Study~16 places
\textsc{State-EXP3} against the strongest published method: ahead in every cell when both are tuned
alike, behind in four of six when that method alone is grid-tuned.

\section{Proof of Theorem~\ref{thm:m1}: additive delay without interleaving}\label{app:m1}

Interleaving cannot give an additive bound. Copy $i$, which absorbs a share $V_i$ of the budget
$V=\sum_iV_i$, has regret $\sqrt{2V_i\log K}$, and
Cauchy--Schwarz over the $m$ copies gives $\sqrt{2mV\log K}$, tight when the $V_i$ agree, so the
$\sqrt m$ in Theorem~\ref{thm:sexp3} is not an artifact of the analysis. An additive bound has to
come from a single copy.

Throughout, $L_t=\sum_{r\le t-d-1}\hat c_r$ and $x_t\propto e^{-\eta L_t}$, so
$L_{t+1}-L_t=\hat c_{t-d}$, with $\hat c_r=0$ for $r<1$. The filtration $\mathcal F_t$ is the
environment's, including outcomes not yet delivered, as in Proposition~\ref{prop:pool}; under it
$x_t$ is $\mathcal F_{t-d-1}$-measurable, so $x_{r+d}$ is $\mathcal F_{r-1}$-measurable.

\paragraph{The drift lemma.} The natural route prices the change of measure from $x_r$ to
$x_{r+d}$ through the normaliser $Z_r$, whose reciprocal has no uniformly bounded moment. That is
avoidable. The pooled estimate satisfies
\[
\langle x_t,\hat c_t\rangle=\sum_ax_t(a)\frac{P(S_t\mid a)X_t}{q_t(S_t)}=X_t\le1
\]
identically, on every path, because the denominator is the state marginal of the same distribution
that supplies the numerator.

\begin{lemma}\label{lem:drift}
If $\eta\le1/(e(d+1))$ then $x_{t+1}(a)\le(1+1/d)\,x_t(a)$ for every $t$ and every $a$, surely, and
hence $x_{t+d}(a)\le e\,x_t(a)$.
\end{lemma}

\begin{proof}
Strong induction on $t$. For $t\le d$ nothing is delivered and $x_{t+1}=x_t$. For $t>d$, assume the
claim for every $s<t$ and compose it over the $d$ steps from $t-d$ to $t$, giving
$x_t(a)\le(1+1/d)^dx_{t-d}(a)\le e\,x_{t-d}(a)$. Contracting with the fixed nonnegative vector
$\hat c_{t-d}$ and using the identity above,
$\langle x_t,\hat c_{t-d}\rangle\le e\langle x_{t-d},\hat c_{t-d}\rangle=eX_{t-d}\le e$. With
$Z_t=\sum_bx_t(b)e^{-\eta\hat c_{t-d}(b)}$ and $e^{-z}\ge1-z$, this gives
$Z_t\ge1-\eta\langle x_t,\hat c_{t-d}\rangle\ge1-\eta e\ge d/(d+1)$, so
$x_{t+1}(a)=x_t(a)e^{-\eta\hat c_{t-d}(a)}/Z_t\le x_t(a)/Z_t\le(1+1/d)x_t(a)$.
\end{proof}

What the induction bounds is
$\langle x_t,\hat c_{t-d}\rangle=X_{t-d}\,q_t(S_{t-d})/q_{t-d}(S_{t-d})$, the state-marginal ratio that the
$1/Z_r$ route could not reach. Individual $\hat c_r(a)$ remain unbounded; only their
$x_t$-weighted mass is controlled. \citet{cesabianchi2019} prove the action-level version under
$\eta\le1/(Ke(d+1))$ and \citet{thune2019} remove the $K$, using that their estimate is carried by
the single action played, so the weight cancels into a same-action ratio. A pooled estimate is
spread over every action and does not collapse that way; the identity above replaces the
cancellation.

\paragraph{Centring.} Let $\zeta_t=\hat c_t-X_t\mathbf 1$, with $\zeta_r=0$ for $r<1$. Exponential
weights is unchanged by a shift that is constant across actions, so $x_t$ is the same iterate and
this is a change of analysis, not of algorithm \citep{eldowa2024}. Keeping
$\gamma(s)=\Ex[X_t^2\mid S_t=s]$ inside the state sum,
\[
\Ex\bigl[\langle x_t,\zeta_t^2\rangle\mid\mathcal F_{t-1}\bigr]
=\sum_s\gamma(s)\Bigl[\frac{\sum_ax_t(a)P(s\mid a)^2}{q_t(s)}-q_t(s)\Bigr]\le v_t-1,
\]
each bracket being nonnegative by Jensen and $\gamma(s)\le1$, the final step being
Remark~\ref{rem:chi2}. Bounding $\langle x_t,\hat c_t^2\rangle$ by $v_t$ first and then subtracting
does not work: it leaves $v_t-\Ex[X_t^2]$.

\begin{proof}[Proof of Theorem~\ref{thm:m1}]
Since $\zeta_r(a)\ge-X_r\ge-1$ we have $\eta\zeta_r(a)\ge-\eta$, and Lagrange's remainder gives
$e^{-z}\le1-z+(e^\eta/2)z^2$ for $z\ge-\eta$. With $\log(1+y)\le y$, the potential
$W_t=\sum_a\exp(-\eta\sum_{r\le t-d-1}\zeta_r(a))$ obeys
$\log(W_{t+1}/W_t)\le-\eta\langle x_t,\zeta_{t-d}\rangle
+(\eta^2e^\eta/2)\langle x_t,\zeta_{t-d}^2\rangle$. Telescoping from $W_1=K$ and using
$\log W_{T+1}\ge-\eta\sum_{r\le T-d}\zeta_r(a^\ast)$ for the comparator $a^\ast$,
\[
\sum_{r=1}^{T-d}\langle x_{r+d},\zeta_r\rangle\le\sum_{r=1}^{T-d}\zeta_r(a^\ast)
+\frac{\log K}{\eta}+\frac{\eta e^\eta}{2}\sum_{r=1}^{T-d}\langle x_{r+d},\zeta_r^2\rangle .
\]
The shift cancels between the two sides. As $x_r$ is $\mathcal F_{r-d-1}$-measurable,
$\Ex[\langle x_r,\hat c_r\rangle]=\Ex[\langle x_r,c_r\rangle]$ and $\Ex[\hat c_r(a^\ast)]=
c_r(a^\ast)$, and the $d$ rounds undelivered by $T+1$ contribute at most $d$, so
\[
\Ex[R_T]\le\frac{\log K}{\eta}
+\frac{\eta e^\eta}{2}\sum_r\Ex\bigl[\langle x_{r+d},\zeta_r^2\rangle\bigr]
+\sum_r\Ex\bigl[\langle x_r-x_{r+d},\hat c_r\rangle\bigr]+d .
\]
Lemma~\ref{lem:drift} gives $x_{r+d}\le e\,x_r$ surely, and $x_{r+d}$ is
$\mathcal F_{r-1}$-measurable, so the second sum is at most $e\sum_r\Ex[v_r-1]\le eV^-$.

For the third, write $x_{r+d}(a)=x_r(a)e^{-\eta G_r(a)}/Z'$ with
$G_r=\sum_{u=r-d}^{r-1}\hat c_u\ge0$ and normaliser $Z'\le1$, distinct from the one-step $Z_t$
of the Lemma. Then
$x_r(a)-x_{r+d}(a)\le x_r(a)(1-e^{-\eta G_r(a)})\le\eta x_r(a)G_r(a)$ for every $a$, trivially
where the left side is negative. Since $\hat c_r\ge0$ and $x_r,G_r$ are
$\mathcal F_{r-1}$-measurable, the conditional expectation is at most $\eta\langle x_r,G_r\rangle$,
and $\Ex[\langle x_r,G_r\rangle]=\sum_{u=r-d}^{r-1}\Ex[\langle x_r,\hat c_u\rangle]\le d$: for each
such $u$, $x_r$ is $\mathcal F_{r-d-1}$-measurable and $r-d-1\le u-1$, so it leaves the conditional
expectation and $\Ex[\langle x_r,\hat c_u\rangle]=\langle x_r,c_u\rangle\le1$. This is
measurability, not independence, and at $u=r-d$ the two $\sigma$-fields coincide, so the step rests
on the timing convention of Section~\ref{sec:setting}. The sum is at most $\eta dT$.

Finally $\eta\le1/(2e)$ gives $e^{1+\eta}<3.3$, so
$\Ex[R_T]\le\log K/\eta+\tfrac{\eta}{2}(3.3V^-+2dT)+d$. For $f(\eta)=A/\eta+B\eta$ restricted to
$\eta\le\eta_{\max}$ one has $\min f\le2\sqrt{AB}+A/\eta_{\max}$, in the interior and in the binding
case alike; with $A=\log K$, $B=(3.3V^-+2dT)/2$ and $\eta_{\max}=1/(e(d+1))$ this is the stated
bound.
\end{proof}

The cap on $\eta$ binds only when $d\gtrsim T/(e^2\log K)$, and not at the settings of
Appendix~\ref{app:num}: the funnel study's own rates are $0.0058$, $0.0031$ and $0.0016$ at
$d=10,50,200$ against caps $0.0334$, $0.0072$ and $0.0018$, so Theorem~\ref{thm:m1} covers those
runs as they were executed. Over three seeds and $20000$ rounds the induction has room to spare,
$\max_ax_{t+1}(a)/x_t(a)$ reaching $1.006$, $1.003$ and $1.002$ against the permitted $1.100$,
$1.020$ and $1.005$.

The two bounds are numerically close on the funnel, where $\bar v\approx2$: the additive form
carries $3.3(\bar v-1)+2d$ and the interleaved one $(d+1)\bar v$, and these cross near $\bar v=2$.
The additive form gains as $\bar v$ rises toward $|\Sset|$, and it describes the algorithm that is
actually run, which is the reason to prefer it.

\end{document}
